\begin{table}[!htbp]
\centering
\caption{Heterogeneous-effect null study: largest deviation (percentage points, with sign) of the rejection probability from the nominal level over the four null values and both scores, the mean estimated variance ratio, and the range of the ratio of the variance of the regenerated statistics to $\widehat V_{{\rm R},n}$.}
\label{tab:hetero-null}
\begin{threeparttable}\begin{singlespace}\small\setlength{\tabcolsep}{4pt}
\begin{tabular}{@{}llrrrrll@{}}
\toprule
Design & Outcome & RT & CRT & SIGA-S & SIGA-R & $\widehat V_{\rm R}/\widehat V_{\rm S}$ & $\widehat V^{*}/\widehat V_{\rm R}$\\
\midrule
2 factors, $n=200$ & continuous & $-0.10$ & $+0.14$ & $-0.14$ & $-0.17$ & 1.003 & 0.998--1.000\\
2 factors, $n=200$ & binary & $-0.74$ & $-0.74$ & $-0.36$ & $-0.37$ & 1.001 & 0.999--1.000\\
2 factors, $n=1000$ & continuous & $-0.15$ & $-0.12$ & $-0.13$ & $-0.15$ & 1.003 & 1.000--1.000\\
2 factors, $n=1000$ & binary & $-0.45$ & $-0.45$ & $-0.13$ & $-0.15$ & 1.001 & 1.000--1.000\\
5 factors, $n=400$ & continuous & $-0.11$ & $+0.12$ & $-0.20$ & $-0.22$ & 1.002 & 0.988--1.001\\
5 factors, $n=400$ & binary & $-0.51$ & $-0.51$ & $-0.20$ & $-0.21$ & 1.001 & 0.997--1.001\\
5 factors, $n=2000$ & continuous & $-0.07$ & $+0.08$ & $+0.08$ & $-0.09$ & 1.002 & 0.997--1.000\\
5 factors, $n=2000$ & binary & $-0.29$ & $-0.28$ & $-0.30$ & $-0.31$ & 1.001 & 0.999--1.000\\
\bottomrule
\end{tabular}
\begin{tablenotes}[flushleft]\small
\item Nominal levels are 5\% (superiority, two-sided; equivalence, two one-sided tests at 5\%) and 2.5\% (non-inferiority); the Monte Carlo standard error of a rejection probability from 100,000 outer trials is 0.07 points at 5\% and 0.05 points at 2.5\%. RT is the reference test and CRT the corrected test with the data-based, noise-adjusted $\widehat{\bm d}(b)$. $\widehat V^{*}$ is the variance of the 4,999 regenerated statistics in a trial, averaged over trials. The negative binary entries for RT and CRT are the unadjusted superiority scenarios, where ties of the lattice-valued score are counted as exceedances. Complete results are in Supplementary Appendix~J.
\end{tablenotes}\end{singlespace}\end{threeparttable}
\end{table}
